% $Id: payload.tex 1784 2012-04-27 23:29:31Z  $
% !TEX root = main.tex

\chapter{CRDT Semantics and the Payload Framework}\label{chap:payload}

This chapter develops the first component of the framework presented in Chapter~\ref{chap:framework}. It formalises the state-based execution model used by CRDTs whose semantics rely solely on semilattice-based state convergence. The definitions introduced here build directly on the algebraic theory developed in Chapter~\ref{chap:theory}, and they serve as the foundation for several of the verified CRDTs presented in Chapter~\ref{chap:implementations}. % Adyust

\section{Introduction}

Conflict-free Replicated Data Types (CRDTs) are data structures that allow multiple replicas to be updated concurrently while guaranteeing eventual convergence. The payload state of a CRDT forms a join-semilattice, and the update operations are monotonic with respect to its partial order. Shapiro \emph{et~al.} formalized this model and showed that if the state merge operation is commutative, associative and idempotent (a join) and the update function monotonically increases the state, then replicas that process the same set of updates (probably in different orders) will reach an identical state~\cite{shapiro2011conflict}. In this chapter we introduce the \texttt{PayloadFramework}, our Athena module that captures these semantics in a reusable form.

\section{Mathematical model}\label{sec:payload-math}

\begin{figure}[t]
  \centering
  \begin{tikzpicture}[
    node distance=1.6cm and 2.2cm,
    box/.style={
      draw,
      rounded corners,
      align=center,
      minimum width=3.5cm,
      minimum height=1.2cm,
      font=\small
    },
    arrow/.style={
      -{Latex[length=3pt]},
      thick
    }
  ]

  % Nodes
  \node[box] (math) {Mathematical model\\[2pt]
    $(S,\sqcup,\leq)$, $\mathsf{apply}$};

  \node[box, right=of math] (abs) {Abstract specification\\[2pt]
    \texttt{PayloadFramework}\\
    (\texttt{merge}, \texttt{leq}, \texttt{apply}, axioms)};

  \node[box, right=of abs] (inst) {Concrete payload\\[2pt]
    \texttt{PayloadOps}\\
    (\texttt{payload-step}, \texttt{payload-join})};

  \node[box, below=of inst] (cfg) {Actor configuration\\[2pt]
    \texttt{Configuration}\\
    (actor states, message net)};

  \node[box, below=of cfg] (runtime) {Runtime behaviour\\[2pt]
    \texttt{Enabled}, \texttt{enqueue}, \texttt{reduce}, \texttt{init}};

  % Arrows (main flow)
  \draw[arrow] (math) -- node[above,align=center,font=\scriptsize]
    {algebraic\\semantics} (abs);

  \draw[arrow] (abs) -- node[above,align=center,font=\scriptsize]
    {instantiation} (inst);

  \draw[arrow] (inst) -- node[right,align=center,font=\scriptsize]
    {payload state\\and updates} (cfg);

  \draw[arrow] (cfg) -- node[right,align=center,font=\scriptsize]
    {message\\processing} (runtime);

  % Feedback arrow (concrete behaviour validates the model)
  \draw[arrow, dashed] (runtime.west) to[bend left=40]
    node[below,align=center,font=\scriptsize]
      {proofs of SEC,\\progress, etc.}
    (abs.south);

  \end{tikzpicture}
  \caption{High-level view of the payload framework. 
  The mathematical model of a join-semilattice with monotone updates
  underpins the abstract \texttt{PayloadFramework} module, which is
  instantiated by concrete payloads (\texttt{PayloadOps}). These are
  plugged into actor configurations and message processing primitives
  to obtain executable CRDT behaviour.}
  \label{fig:payload-overview}
\end{figure}

Let $(S, \sqcup)$ be a join-semilattice: a set $S$ equipped with a binary operation $\sqcup : S \times S \to S$ that is commutative, associative and idempotent. It induces a partial order $\leq$ on $S$ by the rule
\[
  x \leq y \iff x \sqcup y = y.
\]
This makes $(S,\leq)$ a partially ordered set with least upper bounds: for any $x,y\in S$, $x \sqcup y$ is the least element $z$ such that $x \leq z$ and $y \leq z$. A CRDT’s \emph{state} lives in such a semilattice. The \emph{merge} of two replica states corresponds to $\sqcup$; it must therefore be commutative ($x \sqcup y = y \sqcup x$), associative ($x \sqcup (y \sqcup z) = (x \sqcup y) \sqcup z$) and idempotent ($x \sqcup x = x$). Updates are modelled by a function $\mathsf{apply} : S \times U \to S$, where $U$ is the set of update operations. To guarantee convergence, updates must be monotonic: for all $s\in S$ and $u\in U$, we require $s \leq \mathsf{apply}(s,u)$~\cite{shapiro2011conflict}. This ensures that applying the same update multiple times does not revert the state or break convergence.

\section{The \texttt{PayloadFramework} module}\label{sec:payload-code}

The \texttt{PayloadFramework} module declares two domains, \texttt{State} and \texttt{Update}, representing the type of the CRDT’s state and update respectively. It then introduces the merge function \texttt{merge}, the order relation \texttt{leq}, and the update function \texttt{apply}, together with axioms capturing the semilattice and monotonicity requirements.

\AthenaCode{lst:payloadframework-code}{Abstract specification in \texttt{PayloadFramework}}
\begin{lstlisting}[language=Athena,style=athena]
module PayloadFramework {
declare merge: [State, State] -> State
declare leq: [State, State] -> Bool

assert [merge-comm] := (forall x y . merge(x,y) = merge(y,x))
assert [merge-assoc] := (forall x y z .
  merge(x, merge(y,z)) = merge(merge(x,y), z))
assert [merge-idem] := (forall x . merge(x,x) = x)
assert [merge-order] := (forall x y . leq(x,y) <-> merge(x,y) = y)

declare apply: [State, Update] -> State
assert [update-monotone] := (forall s u . leq(s, apply(s,u)))
}
\end{lstlisting}

\noindent\textbf{Semantics.} The definition in Listing~\ref{lst:payloadframework-code} states that \texttt{merge} is a commutative, associative, idempotent operation with \texttt{leq} as the defined order: \texttt{merge(x,y) = y} exactly when \texttt{leq(x,y)}. The axiom \texttt{update-monotone} requires that applying an update never decreases the state. Together, these conditions ensure that the state space $(\texttt{State},\texttt{leq},\texttt{merge})$ forms a join-semilattice and that \texttt{apply} is monotonic. Shapiro \emph{et~al.}~\cite{shapiro2011conflict} proved that any state-based CRDT satisfying these axioms is \emph{strongly eventually consistent}: replicas that see the same set of updates converge to the same state.

\subsection{Instantiating the framework}

The \texttt{Payload} module provides a thin wrapper for instantiating the abstract \texttt{PayloadFramework}. A payload is represented by a structure \texttt{PayloadOps} that carries its state and update types, along with defining equations for the update and merge functions. More precisely, a payload instance \lstinline{(payload-ops init)} stores the initial state \lstinline{init}, and two abstract operations \lstinline{payload-step} and \lstinline{payload-join} must be defined for that payload. The module’s interface can be summarised as shown in Listing~\ref{lst:payloadops-code}:

\AthenaCode{lst:payloadops-code}{Abbreviated interface of the \texttt{Payload} module}
\begin{lstlisting}[language=Athena,style=athena]
module Payload {
  structure (PayloadOps State Value) := (payload-ops State)

  declare payload-step: (State, Value)
    [(PayloadOps State Value) Value State] -> State
  declare payload-join: (State, Value)
    [(PayloadOps State Value) State State] -> State

  define (payload-apply ops value state) :=
    (payload-step ops value state)
}
\end{lstlisting}

\subsection{Payload initialisation}

For ease of use the framework provides helper functions that extract the initial state from a payload descriptor. These definitions do not change the semantics of a payload. Instead, they provide a uniform way to recover the initial state inside generic operations. Listing~\ref{lst:payload-canonical} illustrates how a payload descriptor is initialised:

\AthenaCode{lst:payload-canonical}{Payload initialisation}
\begin{lstlisting}[language=Athena,style=athena]
define (payload-init ops) :=
  match ops {
    (payload-ops init) => init
  }
\end{lstlisting}

\noindent Here \lstinline{payload-init} pattern-matches on the \lstinline{PayloadOps} constant and simply returns the stored initial state. This ensures that every payload descriptor carries the necessary information to initialise actors uniformly. % Adyust

\noindent In this interface, a CRDT is represented by a \lstinline{PayloadOps} constant encapsulating its initial state. A concrete CRDT defines equations for \lstinline{payload-step} and \lstinline{payload-join} by matching that constant. For example the grow-only set defines \lstinline{payload-step} as set insertion and \lstinline{payload-join} as set union. A grow-only counter defines \lstinline{payload-step} as increment and \lstinline{payload-join} as the maximum. The wrapper \lstinline{payload-apply} is a convenience for applying updates during message reduction.

\subsection{Actor configurations and message processing}

The \texttt{Payload} module defines datatypes and functions to build actor configurations with payload states and message networks. A configuration carries a map from actor identifiers to their local states and a multiset of pending messages. By encapsulating these structures and utilities in a reusable module, we can reason about message passing while focusing on the algebraic properties of merge and update operations.

\paragraph{Map and message abstractions.} We first define convenience aliases for empty maps and multisets and overload some operations for finite maps and multisets. We then introduce the \texttt{Message} datatype (a tagged triple consisting of a destination identifier, a payload and some metadata) and the \texttt{Configuration} datatype, which bundles an actor state map and a message network. Two simple selectors extract the actor map and the network from a configuration. These definitions are given in Listing~\ref{lst:payload-details-1}:

\AthenaCode{lst:payload-details-1}{Message and configuration utilities}
\begin{lstlisting}[language=Athena,style=athena]
datatype (Message Value Meta) := (msg Ide Value Meta)

datatype (Configuration State Value Meta) :=
  (kfg am:(FMap.Map Ide State)
       net:(MSet.MSet (Message Value Meta)))

define empty-map := FMap.empty-map
define empty-mset := MSet.null

define (am cfg) :=
  match cfg {
    (Payload.kfg actors _) => actors
  }

define (net cfg) :=
  match cfg {
    (Payload.kfg _ pending) => pending
  }
\end{lstlisting}

The type \lstinline{Message} encapsulates a destination identifier, a payload and optional metadata. A \lstinline{Configuration} pairs an actor state map with a multiset of pending messages. The selector functions \lstinline{am} and \lstinline{net} simply return the state map and the message multiset, respectively. % Adyust

\paragraph{Map lookup.} We next define a helper used internally by the message processing logic. The function \lstinline{map-lookup} in Listing~\ref{lst:payload-map-lookup} searches for a key in a finite map and returns \lstinline{SOME} of the associated value if found, or \lstinline{NONE} otherwise. This function is analogous to the associative array lookup in standard data structures. It is defined by structural recursion on the map, using \lstinline{canonical-map} to normalise the input:

\AthenaCode{lst:payload-map-lookup}{Map lookup helper}
\begin{lstlisting}[language=Athena,style=athena]
define (canonical-map candidate) := candidate

define (map-lookup map target) :=
  match (canonical-map map) {
    FMap.empty-map => NONE
  | (FMap.update (pair key value) rest) =>
      match (equal? key target) {
        true => (SOME value)
      | false => (map-lookup rest target)
      }
  }
\end{lstlisting}

\paragraph{Message reduction.} The core of the execution model is the \lstinline{reduce} function, which delivers a single message and updates the actor state accordingly. If the message network is empty, the configuration is returned unchanged. Otherwise we remove one message \lstinline{(msg id payload meta)} from the multiset, look up the current state of the target actor, and, if the actor exists, apply the payload and update the state map as defined in Listing~\ref{lst:payload-reduce}. Formally:

\AthenaCode{lst:payload-reduce}{Reduction of a single message}
\begin{lstlisting}[language=Athena,style=athena]
define (reduce ops cfg) :=
  match cfg {
    (Payload.kfg actors MSet.null) => cfg
  | (Payload.kfg actors
       (MSet.insert (msg id payload meta) rest)) =>
      match (map-lookup actors id) {
        NONE => (Payload.kfg actors rest)
      | (SOME current) =>
          (Payload.kfg
             (FMap.update
                (pair id (payload-apply ops payload current))
                actors)
             rest)
      }
  }
\end{lstlisting}

\paragraph{Enqueueing new messages.} To schedule updates for delivery, the helpers \lstinline{enqueue-with} and \lstinline{enqueue} add messages to the network. The first allows attaching an arbitrary metadata tag, while the second simply uses the unit tag. These constructors are defined in Listing~\ref{lst:payload-enqueue}:

\AthenaCode{lst:payload-enqueue}{Enqueueing updates for delivery}
\begin{lstlisting}[language=Athena,style=athena]
define (enqueue-with cfg id payload meta) :=
  (Payload.kfg (canonical-map (am cfg))
               (MSet.insert (msg id payload meta)
                            (net cfg)))

define (enqueue cfg id payload) :=
  (enqueue-with cfg id payload unit)
\end{lstlisting}

\paragraph{Initialising actor maps.} Finally, the helper \lstinline{init-map} builds an initial state map given a list of actor identifiers by assigning each actor the payload’s initial state. The \lstinline{init} function combines this map with an empty message set to produce a starting configuration. In Listing~\ref{lst:payload-init} we explicitly declare the type of \lstinline{init} to emphasise that it returns a configuration:

\AthenaCode{lst:payload-init}{Initialising actor maps and configurations}
\begin{lstlisting}[language=Athena,style=athena]
define (init-map ops ids) :=
  match ids {
    nil => empty-map
  | (id :: rest) =>
      (FMap.update
         (pair id (payload-init ops))
         (init-map ops rest))
  | _ => empty-map
  }

declare init: (State, Value)
  [(PayloadOps State Value) (List Ide)]
  -> (Configuration State Value Payload.Unit)

define (init ops ids) :=
  (Payload.kfg (init-map ops ids) empty-mset)
\end{lstlisting}

The helper \lstinline{init-map} traverses the list of actors and uses \lstinline{payload-init} to seed each actor with an initial state. The resulting map is paired with an empty message network by \lstinline{init} to form a configuration ready for execution.

\section{Properties and verification}

The axioms embedded in \texttt{PayloadFramework} guarantee strong eventual consistency for any CRDT that instantiates the module. In particular, it is possible to prove that the merge function is a join operation: it yields the least upper bound of its arguments with respect to \texttt{leq}. We know that by monotonicity, update operations move the state upwards in the lattice. These two facts are sufficient to show that if two replicas process the same multiset of updates (possibly in different orders), their final states are equal. % Adyust

\subsection{Formal Connection to Theory}

The abstract \texttt{PayloadFramework} axioms described in this chapter correspond directly to the formal CRDT theory presented in Chapter~\ref{chap:theory}. The five axioms guarantee Strong Eventual Consistency~\cite{shapiro2011conflict}:

\begin{enumerate}
  \item \texttt{merge-comm} -- Commutativity ensures merge order does not matter.
  \item \texttt{merge-assoc} -- Associativity ensures grouping does not matter.
  \item \texttt{merge-idem} -- Idempotence ensures duplicate merges are harmless.
  \item \texttt{merge-order} -- Defines the partial order from the join operation.
  \item \texttt{update-monotone} -- Ensures updates only move up the lattice.
\end{enumerate}

Concrete payload implementations (\texttt{GCounterPayload}, \texttt{GSetPayload}, etc.) must satisfy these axioms. The theory bridge files in \texttt{src/framework/theory/} provide formal proofs:

\begin{itemize}
  \item \texttt{gcounter\_payload\_theory.ath} proves \texttt{GCounterPayload} satisfies all five axioms using the properties of \texttt{gcounter-max}.
  \item \texttt{gset\_payload\_theory.ath} proves \texttt{GSetPayload} satisfies all five axioms using set union properties.
\end{itemize}

These proofs ensure that the executable implementations are correct by construction and benefit from the SEC theorem proven in the abstract theory layer.

\endinput
